% Propietario conservado: /Users/ruben/Documents/ChatGPT/jueces y controles/REVISION_ARGUMENTAL_APERTURAS_CIERRES_20260915/II/manuscript_es/sections/07c_covarianza_accion.tex
% Propuesta de inserción después de 07b_geometria_elipse.tex.
% RESULTADO_RECUPERADO: c52_06_01_incertidumbre.tex, líneas 10--91.
% Dependencias: c52_04_01_ley_areal.tex, líneas 137--198;
% c52_06_02_intercambio_constitutivo.tex, líneas 5--29.
% Explicitación local: dominio operatorio, prueba de covarianza,
% transporte recíproco y distinción dimensional posición--momento.
% Este fragmento no introduce una constante ni un generador adicional.

\subsection{Canonical covariance and reciprocity of the action ellipse}
\label{ii:sec:ii-area-covarianza}

The preceding construction determines two complementary invariants: the product of the semiaxes fixes the action area and their quotient preserves the anisotropy of the angular pair. We now specify the canonical realization in which these semiaxes are two standard deviations. The starting structure is the ellipse already obtained from the angular and action expressions of APP--TRIT--TPK; the operator realization is established after that construction.

In this subsection we write \(\eta=\eta_{\rm el}
=\operatorname{artanh}(C^*/A)\), with \(|C^*/A|<1\) and \(\hbar>0\). The angular quotient uses the same unit for both its components. In the balanced coordinates \((Q,P)\), both of square-root-of-action dimension, we define
\begin{equation}
 V_\eta:=\frac14
 \begin{pmatrix}a_S^2&0\\0&b_S^2\end{pmatrix}
 =\frac{\hbar}{2}
 \begin{pmatrix}e^\eta&0\\0&e^{-\eta}\end{pmatrix}.
 \label{ii:eq:ii-covarianza-geometrica}
\end{equation}
This positive matrix has entries of action dimension. Its interpretation as a quantum covariance is specified by the following realization.

\begin{proposition}[Gaussian realization of the action covariance]
\label{ii:prop:ii-covarianza-gaussiana}
Consider the canonical representation on \(L^2(\mathbb R,dQ)\), with
\(\widehat Qf(Q)=Qf(Q)\) and
\(\widehat Pf(Q)=-i\hbar f'(Q)\), initially on the Schwartz
space \(\mathcal S(\mathbb R)\). The normalized state
\begin{equation}
 \psi_\eta(Q)=\bigl(\pi\hbar e^\eta\bigr)^{-1/4}
 \exp\!\left(-\frac{Q^2}{2\hbar e^\eta}\right)
 \label{ii:eq:ii-gaussiana-accion}
\end{equation}
has zero means and covariance matrix \(V_\eta\). In particular,
\begin{equation}
 (\Delta Q)^2=\frac{\hbar e^\eta}{2},\qquad
 (\Delta P)^2=\frac{\hbar e^{-\eta}}{2},\qquad
 \operatorname{Cov}(Q,P)=0,\qquad
 \det V_\eta=\frac{\hbar^2}{4}.
 \label{ii:eq:ii-varianzas-accion}
\end{equation}
This realization saturates the Robertson--Schrödinger inequality.
\end{proposition}

\begin{proof}
The preceding operators are essentially self-adjoint on
\(\mathcal S(\mathbb R)\), which is a common invariant core;
on this space,
\([\widehat Q,\widehat P]=i\hbar I\). For
\(\widehat Y=(\widehat Q,\widehat P)\), the symmetric covariance is
\[
 V_{jk}=\frac12\left\langle
 (\widehat Y_j-\langle\widehat Y_j\rangle)
 (\widehat Y_k-\langle\widehat Y_k\rangle)
 +(\widehat Y_k-\langle\widehat Y_k\rangle)
 (\widehat Y_j-\langle\widehat Y_j\rangle)
 \right\rangle.
\]
With \(s=\hbar e^\eta>0\), let \(I_s=\int_{\mathbb R}e^{-Q^2/s}\,dQ\). The product of two equal integrals, evaluated in polar coordinates, gives \(I_s^2=2\pi\int_0^\infty r e^{-r^2/s}\,dr=\pi s\). Moreover, the integral of \(\frac{d}{dQ}(Qe^{-Q^2/s})\) is zero, owing to decay at both endpoints; hence \(\int_{\mathbb R}Q^2e^{-Q^2/s}\,dQ=sI_s/2\). These identities and parity prove normalization, the zero mean and \(\langle\widehat Q^2\rangle=s/2\). Moreover, \(\widehat P\psi_\eta=i e^{-\eta}Q\psi_\eta\). Therefore, \(\langle\widehat P\rangle=0\), \(\|\widehat P\psi_\eta\|^2=\hbar e^{-\eta}/2\) and \(\operatorname{Re}\langle\widehat Q\psi_\eta,
\widehat P\psi_\eta\rangle=0\).

For any normalized state in that domain, Cauchy--Schwarz applied
to \((\widehat Q-\langle\widehat Q\rangle)\psi\) and
\((\widehat P-\langle\widehat P\rangle)\psi\) gives
\[
 (\Delta Q)^2(\Delta P)^2
 \geq \operatorname{Cov}(Q,P)^2
       +\frac14\left|\langle[\widehat Q,\widehat P]\rangle\right|^2
 =\operatorname{Cov}(Q,P)^2+\frac{\hbar^2}{4}.
\]
The computed variances attain equality. Equivalently, the
operator
\[
 a_\eta=\frac{e^{-\eta/2}\widehat Q
                  +i e^{\eta/2}\widehat P}{\sqrt{2\hbar}}
\]
satisfies \([a_\eta,a_\eta^\dagger]=I\) and
\(a_\eta\psi_\eta=0\), by direct substitution.
\end{proof}

\begin{proposition}[Area, dispersion, and reciprocal transport]
\label{ii:prop:ii-area-covarianza-reciprocidad}
The action ellipse is exactly the covariance sublevel set
\begin{equation}
 \mathcal E_\eta
 =\left\{y=(Q,P)^{\mathsf T}:y^{\mathsf T}V_\eta^{-1}y\leq4\right\}
 =\left\{\frac{Q^2}{a_S^2}+\frac{P^2}{b_S^2}\leq1\right\}.
 \label{ii:eq:ii-elipse-covarianza}
\end{equation}
Its area is \(h=2\pi\hbar\), its semiaxes are
\(a_S=2\Delta Q\), \(b_S=2\Delta P\), and
\begin{equation}
 \frac{\Delta P}{\Delta Q}=e^{-\eta}
 =\frac{b_S}{a_S}=R_{\rm el}^2.
 \label{ii:eq:ii-covarianza-forma}
\end{equation}
The canonical rotation \(J_2\) of Section~\ref{ii:sec:ellipse} transports
\(V_\eta\) to \(V_{-\eta}\) and preserves the symplectic form.
\end{proposition}

\begin{proof}
Inverting the diagonal matrix
\eqref{ii:eq:ii-covarianza-geometrica} gives
\eqref{ii:eq:ii-elipse-covarianza}. Its semiaxes are the positive roots of
four times its variances, and
\[
 \operatorname{Area}(\mathcal E_\eta)
 =4\pi\sqrt{\det V_\eta}=2\pi\hbar=h,
\]
in agreement with the areal law~\eqref{ii:eq:ii-elipse-area-fase}.
The ratio of the positive roots proves
\eqref{ii:eq:ii-covarianza-forma}.
If \(M_\eta=\operatorname{diag}(e^{\eta/2},e^{-\eta/2})\), then
\[
 V_\eta=M_\eta V_0M_\eta^{\mathsf T},\qquad
 M_\eta^{\mathsf T}J_2M_\eta=J_2,\qquad
 J_2V_\eta J_2^{\mathsf T}=V_{-\eta}.
\]
These identities follow by multiplication. In operator terms,
\((\widehat Q',\widehat P')=(-\widehat P,\widehat Q)\) preserves
\([\widehat Q',\widehat P']=i\hbar I\). The permutation
\(S_2:(Q,P)\mapsto(P,Q)\) also exchanges the variances but reverses
the sign of the commutator and of the symplectic form. Thus covariance
transport respects the distinction between canonical reciprocity and
orientation reversal already established in Section~\ref{ii:sec:ellipse}.
\end{proof}

\paragraph{Position, momentum, and dimensional transduction.}
Let \(\lambda_{xp}>0\) be a factor of the dimensional chart, with the dimension of momentum per unit length. The transformation
\[
 Q=\sqrt{\lambda_{xp}}\,x,\qquad
 P=\frac{p_x}{\sqrt{\lambda_{xp}}}
\]
realizes the balanced pair as a position \(x\) and its conjugate momentum
\(p_x\); it preserves \(dQ\wedge dP=dx\wedge dp_x\) and
\([\widehat x,\widehat p_x]=i\hbar I\). By transformation of
covariance,
\begin{equation}
 V_\eta^{(x,p_x)}=\frac{\hbar}{2}
 \begin{pmatrix}e^\eta/\lambda_{xp}&0\\
                 0&\lambda_{xp}e^{-\eta}\end{pmatrix},\qquad
 \Delta x\,\Delta p_x=\frac{\hbar}{2},\qquad
 \frac{\Delta p_x}{\Delta x}=\lambda_{xp}e^{-\eta}.
 \label{ii:eq:ii-covarianza-dimensiones}
\end{equation}
The dispersion ratio acquires the dimension of momentum per length here. By contrast, \eqref{ii:eq:ii-covarianza-forma} belongs to the balanced chart. For the inductive--capacitive realization already defined, \eqref{ii:eq:ii-elipse-impedancia} supplies the dimensionless identity \(Z_\eta/Z_*=e^{-\eta}=\Delta P/\Delta Q\). Each reading uses its declared transduction and preserves the same action determinant.

\paragraph{Contour energy and mean energy of the state.}
The inductive--capacitive realization has Hamiltonian
\[
 H_\eta(Q,P)=\frac{\omega}{2}
 \left(e^{-\eta}Q^2+e^\eta P^2\right),\qquad \omega>0.
\]
On \(\partial\mathcal E_\eta\), its value is
\(H_\eta=\hbar\omega\), by the semiaxis equations. In the
quantum realization of Proposition~\ref{ii:prop:ii-covarianza-gaussiana},
\begin{equation}
 \langle\widehat H_\eta\rangle_{\psi_\eta}
 =\frac{\omega}{2}
 \left(e^{-\eta}\frac{\hbar e^\eta}{2}
       +e^\eta\frac{\hbar e^{-\eta}}{2}\right)
 =\frac{\hbar\omega}{2}.
 \label{ii:eq:ii-energia-contorno-covarianza}
\end{equation}
The two-standard-deviation contour therefore has twice the
mean energy of this Gaussian state. The areal law fixes action
normalization; the angular pair fixes its distribution between the two coordinates;
the canonical representation and explicit state realize that distribution
as covariance. Saturation of uncertainty is thus established in
the realization specified by its operators and its state.
